\documentclass[10pt,a4paper]{article}
\usepackage[margin=20mm]{geometry}
\usepackage[T1]{fontenc}
\usepackage{lmodern,amsmath,amssymb,booktabs,array,tabularx,longtable,xcolor}
\usepackage{enumitem,microtype,fancyhdr,xurl}
\usepackage[colorlinks=true,linkcolor=blue!45!black,urlcolor=blue!45!black]{hyperref}
\hypersetup{pdftitle={Statement audit: body.tex and the Lean formalization},pdfauthor={Codex}}
\definecolor{ink}{HTML}{18344B}
\definecolor{accent}{HTML}{236B65}
\pagestyle{fancy}\fancyhf{}
\fancyhead[L]{\small\color{ink}Lanczos statement audit}
\fancyhead[R]{\small 18 October 2026}
\fancyfoot[L]{\footnotesize Local source snapshot; references are source line numbers}
\fancyfoot[R]{\thepage}
\setlength{\headheight}{14pt}
\setlength{\parindent}{0pt}
\setlength{\parskip}{6pt}
\setlength{\emergencystretch}{2em}
\setlist[itemize]{leftmargin=*,itemsep=3pt,topsep=3pt}
\newcommand{\code}[1]{{\small\nolinkurl{#1}}}
\newcommand{\norm}[1]{\left\lVert#1\right\rVert}
\newcommand{\eps}{\varepsilon}
\newcommand{\loc}[2]{\code{#1}, line #2}
\newcommand{\head}[1]{\vspace{4pt}{\Large\bfseries\color{ink}#1}\par\vspace{3pt}}
\newcommand{\subhead}[1]{\vspace{3pt}{\bfseries\color{accent}#1}\par}
\newcolumntype{Y}{>{\raggedright\arraybackslash}X}
\begin{document}

{\huge\bfseries\color{ink}Do the paper's claims match Lean?}\par
{\large Statement audit of \code{body.tex}}\par
18 October 2026\quad\textbullet\quad 42 Lean source files, 6,632 lines

\head{Verdict}
The main mathematical results are substantially aligned. Each of the paper's
\textbf{13 labelled theorems, lemmas, and corollaries} has a corresponding Lean
result, with explicit constants for the bounds hidden by $\lesssim_k$.
The backward construction, exact Lanczos output, localization, containment,
and the Lipschitz spectral-measure estimate are supported by actual proofs.

\textbf{The paper does not fully match as written.} Its divided Ritz-residual
inequality omits a nonzero-vector hypothesis that the Lean theorem explicitly
requires. A permitted run can have $V_k y=0$ for a unit eigenvector $y$ of $T_k$.
The general commutator bound also needs an empty-maximum convention at $k=1$.
The claim that Lean proves ``all of the statements in this paper'' is too broad
for this source tree.

\subhead{Verification evidence}
\begin{itemize}
\item \texttt{lake build} succeeded: \textbf{3,255 jobs}, with linter warnings.
The build used the existing local dependency/build cache.
\item An executable Lean audit checked the transitive axiom dependencies of
\textbf{1,156 declarations} under \code{FinitePrecisionLanczos}, including
\textbf{890 theorem declarations} (these counts include generated declarations).
Only \code{propext}, \code{Classical.choice}, and \code{Quot.sound} occurred.
\item A source scan of all 42 project Lean files found no occurrences of
\code{sorry}, \code{admit}, \code{axiom}, \code{unsafe},
\code{implemented_by}, or \code{native_decide} as whole words.
\item The root import reaches every project module. The backward-model import
closure does not include \code{LossOfOrthogonality}; the Paige proof's import
closure does not include the normalized-vector modules.
\end{itemize}

\subhead{Findings at a glance}
\begin{tabularx}{\textwidth}{@{}p{17mm}Yp{27mm}@{}}
\toprule
Priority & Finding & Paper location\\\midrule
Required & Add $V_k y\ne0$ before dividing by $\norm{V_k y}$. & Lines 585--594\\
Required & Restrict the blanket formalization claim to the statements actually covered. & Line 102\\
Minor & Define $\max\varnothing=0$ for the nonnegative defect-jump maximum. & Lines 510--515\\
Clarify & Distinguish Ritz-value localization from normalized Ritz-vector convergence; quantify substantial loss of orthogonality. & Line 596\\
Scope & Lean assumes a run satisfying a kernel-error model; it does not verify an IEEE implementation or the experiments. & Lines 108--259; 1043--1062\\
\bottomrule
\end{tabularx}

\subhead{What this report compares}
The supplied \code{body.tex} and \code{lean/FinitePrecisionLanczos*.lean}
source tree are authoritative for this audit. The README was used as a navigation
aid, not as proof evidence. The paper and Lean sources were left unchanged.
Historical attributions and cited external papers were not independently audited.
The referenced \code{AI} include and experiment assets are absent from this
workspace snapshot.

\newpage
\head{1. A real missing hypothesis}
\subhead{Divided residual: paper lines 585--594}
For $x=V_k y$, the paper displays
\[
 \frac{\norm{(A-\theta I)x}}{\norm{x}}
 \le \frac{\sqrt{1+\eps}\,\beta_k|e_k^T y|+\norm{F_k}}{\norm{x}}
\]
without assuming $x\ne0$. In \loc{Lanczos/Ritz.lean}{146},
\code{ritz_residual_bound_div} has the additional hypothesis
\begin{quote}\verb|(hx : 0 < vnorm (L.ritzVector y))|.\end{quote}
Its undivided counterpart, \code{ritz_residual_bound} (line 126), needs no such
hypothesis. The condition cannot be inferred merely from $\norm{y}=1$:
the computed columns may be linearly dependent.

\subhead{Explicit admissible example}
Take $n=1$, $k=k_{\max}=2$, $A=[1]$, and $\eta=\eps=1/1000$.
All vector quantities below are scalars. Set
\[
 v_0=0,\quad v_1=v_2=1,\quad v_3=0,\qquad
 \beta_0=0,\quad\beta_1=\eta,\quad\beta_2=0.
\]
The two steps are
\[
\begin{array}{c|ccccc}
j&u_j&w_j&\alpha_j&z_j&e^\alpha_j\\\hline
1&1&1&1-\eta&\eta&-\eta\\
2&1&1-\eta&1-\eta&0&0
\end{array}
\]
All other arithmetic errors and normalization errors are zero. These data satisfy
the kernel identities and bounds: in particular
$|e^\alpha_1|=\eps\norm{v_1}\norm{w_1}$,
$\beta_1v_2=z_1$, and $2\eps\le1/100$.
The final step is terminal, exactly as allowed by the paper.
But
\[
 T_2=\begin{pmatrix}1-\eta&\eta\\\eta&1-\eta\end{pmatrix},\qquad
 y=\frac1{\sqrt2}\binom{1}{-1},\quad
 \theta=1-2\eta,\qquad V_2y=0.
\]
Thus the displayed ratios and $x/\norm{x}$ are undefined in ordinary real analysis.
The example's model identities were independently checked with exact rational
arithmetic in \code{check_zero_ritz.py}; it is not a claim about which errors a
particular IEEE implementation would produce. Lean's totalized division at zero
would not repair the mathematical interpretation.

\subhead{Recommended change}
Before the divided display, insert: ``If $x\ne0$, dividing the undivided residual
bound by $\norm{x}$ gives\ldots'' Keep the undivided inequality as the unconditional
statement. Apply the same qualification to the normalized-vector discussion.

\subhead{The convergence sentence needs care, not a new main theorem}
Line 596 identifies a small scalar residual with convergence of the Ritz pair.
The proved Paige result is the product bound
$s\rho\le C_k\eps\norm{A}$, where
$s=|x^Tv_{k+1}|$, $\rho=\beta_k|e_k^Ty|$, and $C_k=3000k^3$.
A small $\rho$ alone does not control the divided residual.
For a nonterminal step and a specified $s\ge\tau>0$, however,
$\norm{x}\ge\tau/\sqrt{1+\eps}$, giving
\[
\frac{\norm{(A-\theta I)x}}{\norm{x}}
\le \frac{(1+\eps)C_k\eps\norm{A}}{\tau^2}
   +\frac{\sqrt{1+\eps}\norm{F_k}}{\tau}.
\]
This is a mathematical consequence of the checked bounds, not a separate theorem
in this project. Alternatively, cite \code{thm:stabilized} for the conclusion that
the \emph{Ritz value} is near the spectrum, which requires no lower bound on $\norm{x}$.

\newpage
\head{2. Assumptions and boundary cases}
\subhead{The common model is consistent}
\texttt{LanczosRun n k}, defined in \loc{Lanczos/Run.lean}{9}, stores the matrix,
vectors, coefficients, local errors, and their hypotheses. Paper $A$ is Lean
\code{L.H}; paper $\eps$ is \texttt{L.}$\varepsilon$. The vector norms use Euclidean space,
and the matrix norms use \code{Matrix.Norms.L2Operator}; they are the paper's
induced 2-norms, not entrywise norms.

\begin{tabularx}{\textwidth}{@{}p{53mm}YY@{}}
\toprule
Paper hypothesis & Lean interface & Assessment\\\midrule
$A=A^T$; real data & \code{hHsymm}; real matrices and vectors & Matches.\\
$0\le\eps\le1/(100k_{\max})$, completed prefix $k\le k_{\max}$ & \texttt{h}$\varepsilon$\texttt{\_nonneg}, \code{hk_pos}, \code{hsmall}: $k\eps\le1/100$ & Paper implies Lean's prefix condition for $k\ge1$.\\
Kernel equations and bounds & \code{h_u}, \code{heu}, \code{h_w}, \code{hew}, \code{h_al}, \code{heal}, \code{h_z}, \code{hez} & Matches; $|\beta|$ in Lean equals $\beta$ by nonnegativity.\\
Normalization equation and error & \code{h_be}, \texttt{h}$\Delta$ & Matches.\\
Near-unit formed columns; conditional last column & \code{hnear}, \code{hnext} & Matches, including terminal case.\\
$\beta_0=0$; internal $\beta_j>0$; final $\beta_k\ge0$ & \texttt{h}$\beta$\texttt{\_zero}, \texttt{h}$\beta$\texttt{\_internal}, \texttt{h}$\beta$\texttt{\_nonneg} & Matches.\\
Initial normalization of original input $v$ & No original-input field or normalization routine & Only the near-unit $v_1$ property is represented.\\
\bottomrule
\end{tabularx}

\subhead{Termination is handled, with a choice of auxiliary vector}
Paper lines 195--259 allow $\beta_k=0$ and set the unformed $v_{k+1}=0$.
Lean does not require that terminal raw vector to be zero; it is used multiplied
by $\beta_k$ where needed. This weakens the assumptions and causes no conflict
for the paper's run.

For the backward construction the paper allows any unit $\bar v_{k+1}$ and sets
$d_{k+1}=1$ (line 677). Lean makes the specific permitted choice
$\bar v_{k+1}=v_1/\norm{v_1}$ in \code{normalizedVec}
(\loc{Lanczos/Normalized/Normalization.lean}{138}). The packaged run theorems
therefore certify that choice. The generic dilation algebra is parameterized by
matrices and supports the paper's construction, but there is no packaged
run-level theorem quantifying over every terminal unit-vector choice.
For literal correspondence, specify the same choice in the paper or generalize
the run-level interface.

\code{GreenbaumBackwardModel.exact_next} deliberately assumes $\beta_k>0$.
The theorem unconditionally reproduces $T_k$, $\beta_k$, and the first $k$ exact
vectors; it does not identify an auxiliary unit vector with an exact Lanczos
output after termination. This agrees with the stated output claim in
\code{thm:greenbaum}. The final sentence of its proof (line 954) should be read
with this distinction.

\subhead{Small indexing clarification in the generic commutator lemma}
At $k=1$, $\max_{1\le i<k}\beta_i|\gamma_{i+1}-\gamma_i|$ is a maximum over
an empty set. The paper gives no convention. Lean's \code{defectJump}
(\loc{Core/SaadGramCommutator.lean}{40}) inserts zero at the final index and
takes a maximum over a nonempty \texttt{Fin k}. Consequently this term is zero
when $k=1$. State that convention, and write the lower index limits explicitly.
Lean uses $|T_{i,i+1}(\gamma_{i+1}-\gamma_i)|$; this equals the printed
expression because the paper's internal $\beta_i$ are positive. If the lemma is
restated for arbitrary signed off-diagonal entries, keep the absolute value on
the coefficient. Also make $1\le k\le k_{\max}$ explicit in the setup.

\newpage
\head{3. Correspondence: local analysis and Paige}
All paths below are relative to \code{lean/FinitePrecisionLanczos/}.
Names in the run-specific rows belong to \code{FinitePrecisionLanczos.LanczosRun}.
``Matches'' includes substituting the explicit constants into the paper's
$\lesssim_k$ convention, and the assumptions described on page 3.

\begin{longtable}{@{}p{44mm}p{121.7mm}@{}}
\toprule
Paper result & Checked Lean statement and assessment\\\midrule
\code{lem:coefficients}\newline lines 289--417 &
\code{coefficients_paper}, \loc{Lanczos/Bounds.lean}{449}, proves
\[
\norm{u_j}\le2\norm A,\quad\norm{w_j}\le5j\norm A,\quad
|\alpha_j|\le6j\norm A,
\]
\[
\norm{z_j}\le7j\norm A,\qquad\beta_j\le4j\norm A.
\]
These match the detailed proof as well as the weaker lemma statement.
\code{beta_recurrence_paper} (line 389) proves
$\beta_{j+1}\le(1+10\eps)\beta_j+2\norm A$.
\code{coefficient_multipliers} (line 10) proves the two multiplier
inequalities on the larger interval $[0,1/20]$. \textbf{Matches.}\\\addlinespace

\code{thm:local-errors}\newline lines 421--468 &
\code{local_errors}, \loc{Lanczos/Bounds.lean}{714}, proves for $1\le j\le k$
\[
\norm{f_j}\le700j\eps\norm A,\quad |g_j|\le\eps,\quad
|p_j|\le1200j\eps\norm A.
\]
The error-sum and cancellation identities are proved in
\code{recurrence_error_identity} and \code{local_product_identity}
(\code{Lanczos/Run.lean}, lines 72 and 93). \textbf{Matches.}\\\addlinespace

\code{lem:gram-commutator}\newline lines 502--565 &
\code{SaadGramRecurrence.gram_commutator},
\loc{Core/SaadGramCommutator.lean}{141}, proves upper triangularity and
$\norm N\le2k(P+c\norm Z+D)$, where $P$ is the maximum adjacent product,
$c$ is the maximum norm of the first $k$ columns, and $D$ is the
maximum absolute defect jump, padded by zero. The rectangular identity is
\code{padded_commutator} (line 103). Symmetry, tridiagonality, Gram splitting,
and the recurrence are explicit structure hypotheses. \textbf{Matches after
the $k=1$ convention on page 3.}\\\addlinespace

\code{cor:gram-commutator}\newline lines 495--498 &
\code{computed_gram_commutator},
\loc{Lanczos/ComputedCommutator.lean}{178}, supplies the printed matrix
definition, the Gram split, upper triangularity, and
$\norm{M_k}\le3000k^3\eps\norm A$.
\textbf{Matches.}\\\addlinespace

\code{thm:paige}\newline lines 612--633 &
\code{paige_bound}, \loc{Lanczos/LossOfOrthogonality.lean}{41}, proves
$|\langle V_ky,\beta_kv_{k+1}\rangle y_k|
\le3000k^3\eps\norm A$ for a unit eigenpair.
Nonnegative $\beta_k$ makes this exactly the paper's product.
\code{computed_paige_identity} (line 26) gives the signed identity
$\langle V_ky,\beta_kv_{k+1}\rangle y_k=-y^TM_ky$.
\textbf{Matches.}\\\addlinespace

\code{lem:normalized-errors}\newline lines 694--716 &
\code{normalized_local_errors},
\loc{Lanczos/Normalized/LocalErrors.lean}{144}, proves
$\norm{\bar f_j}\le1500k\eps\norm A$ and
$|\bar p_j|\le4800k\eps\norm A$.
The same computed $\alpha_j,\beta_j$ are used throughout; there is no
replacement boundary coefficient. \textbf{Matches for the chosen terminal
completion.}\\\addlinespace

\code{cor:normalized-commutator}\newline lines 732--740 &
\code{normalized_gram_commutator},
\loc{Lanczos/Normalized/Commutator.lean}{104}, proves upper triangularity and
$\norm{\bar M_k}\le9600k^3\eps\norm A$ using the normalized recurrence and
unit columns. \textbf{Matches for the chosen terminal completion.}\\
\bottomrule
\end{longtable}

\newpage
\head{4. Correspondence: backward model and Ritz values}
\begin{longtable}{@{}p{44mm}p{121.7mm}@{}}
\toprule
Paper result & Checked Lean statement and assessment\\\midrule
\code{lem:dilation}\newline lines 767--809 &
\code{physical_dilation}, \loc{Lanczos/Greenbaum/Statements.lean}{23},
proves $\widehat V_{k+1}^T\widehat V_{k+1}=I$ and the prescribed first column.
The definitions in \code{Greenbaum/Dilation.lean} use exactly $k+1$ blocks
$\bar V_{k+1}C_{k+1}K_{k+1}^r$, $0\le r\le k$.
\code{triangularC} is a finite geometric sum, with both inverse identities
proved in \code{Core/Nilpotent.lean}; this is the paper's inverse, not a
different construction. \textbf{Matches for the chosen terminal completion.}\\\addlinespace

\code{lem:exchange}\newline lines 824--864 &
\code{exchange_rule}, \loc{Lanczos/Greenbaum/Statements.lean}{38}, proves
$\norm{K_k}\le1$, $\norm{C_k}\le2$,
$\norm{L_k}\le38400k^3\eps\norm A$, and for each $r\in\mathbb N$ an exchange
matrix $D_r$ with $\norm{D_r}\le r\norm{L_k}$.
\code{greenbaumD_zero} and \code{greenbaumD_succ}
(\code{Greenbaum/Dilation.lean}, lines 201 and 203) prove the printed recurrence
for $D_r$. \textbf{Matches.}\\\addlinespace

\code{lem:dilated-recurrence}\newline lines 868--904 &
\code{dilated_recurrence}, \loc{Lanczos/Greenbaum/Statements.lean}{76}, proves
$\widehat A\widehat V_k=\widehat V_{k+1}\underline T_k+\widehat E$ with
$\norm{\widehat E}\le240000k^6\eps\norm A$.
\code{greenbaumErrorBlock_eq} gives precisely
$E_r=\bar F_k C_kK_k^r+\bar V_k(D_r-D_{r+1})$.
The sharper $\sqrt k$ and $\sqrt{k+1}$ intermediate estimates are also present
in \code{greenbaumErrorBlock_norm_sqrt} and
\code{greenbaumError_norm_sqrt} (\code{Greenbaum/Dilation.lean}, lines 241 and
251). \textbf{Matches.}\\\addlinespace

\code{thm:greenbaum}\newline lines 911--955 &
\code{greenbaum_backward_model}, \loc{Lanczos/Greenbaum/Backward.lean}{43},
proves dimension $n(k+1)$, symmetry, the corrected recurrence,
orthonormality, the starting vector, exact output $T_k$ and $\beta_k$, and
no breakdown before step $k$.
The perturbation satisfies $\norm\Delta\le480000k^6\eps\norm A$.
\code{greenbaumDelta_eq} and \code{greenbaumDelta_projector}
(\code{Greenbaum/Correction.lean}, lines 66 and 54) prove both of the paper's
formulas, with $\Delta=\widetilde A-\widehat A$.
\code{greenbaumDelta_norm} proves $\norm\Delta\le2\norm{\widehat E}$.
\textbf{Matches; terminal output is qualified as on page 3.}\\\addlinespace

\code{thm:stabilized}\newline lines 981--1008 &
\code{greenbaum_localization_exists},
\loc{Lanczos/Greenbaum/Corollaries.lean}{86}, proves that an actual eigenvalue
$\lambda$ of $A$ satisfies
\[
 |\theta-\lambda|\le\beta_k|y_k|+480000k^6\eps\norm A.
\]
The sharper perturbation form and the exact dilated Ritz-residual identity
also occur in this file. This is at least as strong as the printed
$\lesssim_k$ bound on the sum. \textbf{Matches.}\\\addlinespace

\code{thm:containment}\newline lines 1013--1039 &
\code{greenbaum_containment_eigenvalue},
\loc{Lanczos/Greenbaum/Corollaries.lean}{161}, proves, for every eigenvalue
$\theta$ of $T_k$,
\[
\lambda_{\min}(A)-480000k^6\eps\norm A\le\theta
\le\lambda_{\max}(A)+480000k^6\eps\norm A.
\]
The code constructs a unit eigenvector for each enumerated eigenvalue;
the conclusion is not limited to eigenpairs supplied externally.
\textbf{Matches.}\\
\bottomrule
\end{longtable}

\newpage
\head{5. The spectral-measure remark is supported}
\code{rem:wasserstein}, paper lines 958--964, contains several distinct claims.
The code proves their substantive analytic content, including the actual matrix
functional calculus and the rank bound.

\subhead{Eigenvalue localization of the backward matrix}
\code{greenbaumTildeA_eigenvalue_localization}
(\loc{Lanczos/Greenbaum/Spectral.lean}{37}) proves
\[
 \operatorname{dist}(\widetilde\lambda,\operatorname{spec}(A))
 \le\norm\Delta
\]
for every eigenvalue of $\widetilde A$. This matches the paper's one-sided
localization conclusion. The proof uses a symmetric residual estimate, rather
than a separately invoked ordered-eigenvalue version of Weyl's theorem.
It does not package a multiplicity-preserving matching of all eigenvalues into
clusters, which is a different statement from this one-sided bound.

\subhead{Equality of the unperturbed starting spectral measures}
\code{starting_spectralMass}
(\loc{Lanczos/Greenbaum/SpectralMass.lean}{14}) proves the equality of the
finite atomic spectral masses of $(\widehat A,\widetilde v)$ and
$(A,\bar v_1)$. The mass at an eigenvalue is the sum of the squared starting-vector
coordinates over its eigenspace, so repeated eigenvalues are combined correctly.
\code{starting_functional_calculus}
(\loc{Lanczos/Greenbaum/Wasserstein.lean}{28}) then proves, for every real function
$f$,
\[
 \bar v_1^T f(A)\bar v_1=\widetilde v^T f(\widehat A)\widetilde v.
\]
Here $f(A)$ is Mathlib's matrix \code{cfc}, with its eigenbasis interpretation
proved in \code{Core/Measure/FunctionalCalculus.lean}.

\subhead{The full inequality chain}
For every $1$-Lipschitz $f$, the formalization proves
\begin{align*}
\big|\bar v_1^Tf(A)\bar v_1-\widetilde v^Tf(\widetilde A)\widetilde v\big|
&\le\norm{f(\widehat A)-f(\widetilde A)}_{\mathrm F}\\
&\le\norm\Delta_{\mathrm F}
\le\sqrt{2k}\,\norm\Delta.
\end{align*}
The relevant declarations are:
\begin{itemize}
\item \code{cfc_lipschitz_frobNorm},
\loc{Core/Measure/FunctionalCalculus.lean}{82}: the full Frobenius-norm
Lipschitz inequality for real symmetric matrices.
\item \code{greenbaumDelta_rank},
\loc{Lanczos/Greenbaum/Correction.lean}{74}: $\operatorname{rank}\Delta\le2k$.
\item \code{frobNorm_le_sqrt_rank}, \loc{Core/SpectralFrobenius.lean}{28}:
the symmetric rank-to-Frobenius estimate.
\item \code{starting_lipschitz_frobNorm} and \code{starting_lipschitz},
\loc{Lanczos/Greenbaum/Wasserstein.lean}{37} and line 48:
the exact scalar bounds printed in the paper.
\end{itemize}

\subhead{Meaning of ``Wasserstein distance'' in this project}
\code{atomicWasserstein} (\loc{Core/Measure/Wasserstein.lean}{26}) is defined as
the supremum of absolute integral differences over all $1$-Lipschitz functions.
\code{starting_wasserstein} proves its bound by $\sqrt{2k}\norm\Delta$;
\code{starting_wasserstein_explicit} inserts $480000k^6\eps\norm A$.
This matches the paper's explicitly stated Lipschitz-test formulation.

There is no project theorem identifying this finite-atomic dual definition with
an infimum over transport couplings or a library probability-measure Wasserstein
metric. That is a formalization boundary, not a failure of the displayed
Lipschitz inequality. If that identification is intended to be part of the
formalization claim, state the dual convention or add the bridge theorem.

\newpage
\head{6. Coverage limits and recommended wording}
\subhead{The ``all statements'' claim is not supported literally}
Paper line 102 says that formal Lean proofs are supplied for all statements.
The source tree supports the numbered analytic results described in this report,
subject to the qualifications above, but does not contain all of the following:
\begin{itemize}
\item A formal execution semantics for Algorithm 1, an IEEE arithmetic model,
or a theorem constructing \code{LanczosRun} from an execution of that algorithm.
The kernel bounds and normalization properties are hypotheses, as the paper's
assumption-based analysis intends.
\item A theorem relating the given original starting vector to the first
computed vector through the initial normalization perturbation (lines 190--193).
The run model starts with its stored near-unit $v_1$.
\item A general theorem that exact Lanczos on an arbitrary symmetric input
produces an orthonormal Krylov basis, including the span equality in lines
24--30. \code{Core/ExactLanczos.lean} defines a process and identifies its outputs
\emph{given} an orthonormal recurrence; its structure
\code{OrthonormalRecurrence} assumes orthogonality. This is sufficient for the
backward theorem, where orthogonality is independently proved by dilation.
\item Verification of the numerical values, cluster counts, or weight errors
in the figures and numerical illustration (lines 41--44, 603--608, 1043--1062).
There is no instantiation of the Lean model with those experiment data.
\item Proofs of historical comparisons, claims about practical accuracy,
or statements about the precision needed for kernels and downstream algorithms.
Those require their own evidence; they are not consequences certified by the
current project merely because they occur in the paper.
\end{itemize}

Suggested replacement for line 102:
\begin{quote}
The accompanying Lean development proves counterparts of the main local-error,
Gram-commutator, loss-of-orthogonality, backward-model, localization, containment,
and Lipschitz spectral-measure results, with explicit constants, for runs
satisfying the stated arithmetic model. The implementation-to-model connection,
introductory background results, and numerical experiments are not formalized.
\end{quote}

\subhead{Numerical illustrations need separate validation}
The formal perturbation theorem is about the exact real matrices in its
construction. Forming the model in double precision, as described in the paper,
introduces additional errors that the Lean theorem does not bound. Likewise,
double precision with full reorthogonalization is a numerical reference for
exact Lanczos, not an exact-arithmetic certificate (figure caption, line 42).
The observed $4\cdot10^{-6}$ weight agreement and $2.5\cdot10^{-7}$ perturbation
norm should be presented as experimental measurements.

The one-sided localization theorem does not by itself certify that each
original eigenvalue has a cluster of exactly $k+1$ numerical model eigenvalues.
Such a count needs a matching/separation argument or direct validated numerical
evidence. This audit makes no claim that the reported experiment is incorrect.

\subhead{Documentation drift}
\code{lean/README.md}, lines 30--31, says the nonzero condition on the divided
Ritz residual appears in both paper and code. It is present in the code but
absent from this \code{body.tex}. The README also links to a correspondence map
and prior reports not supplied in this workspace. Those assertions were not
used to establish the verdict.

\subhead{Minimum changes to align the paper}
Add the nonzero condition, define the empty defect-jump maximum, narrow the
formalization claim, and qualify the convergence sentence. For exact statement
correspondence at termination, specify the same auxiliary unit-vector choice as
Lean. None of these changes requires weakening the proved main stability bounds.

\newpage
\head{7. Traceability and reproduction}
\subhead{Additional displayed identities checked}
\begin{tabularx}{\textwidth}{@{}p{53mm}Y@{}}
\toprule
Paper equation or proof step & Lean evidence (paths relative to the project namespace)\\\midrule
\code{eq:governing} & \code{governing}, \code{Lanczos/Ritz.lean}:63\\
\code{eq:gram-split}, \code{eq:M-def} & \code{computed_gram_split}, \code{computedCommutator}, \code{Lanczos/ComputedCommutator.lean}:31,63; generic \code{padded_commutator}\\
\code{eq:gram-commutator-full} & \code{GramRecurrence.gram_commutator_full}, \code{Core/GramCommutator.lean}:129\\
\code{eq:generic-commutator} & \code{SaadGramRecurrence.N_skew}, \code{Core/SaadGramCommutator.lean}:112\\
\code{eq:ritz-residual-identity} & \code{ritz_residual_identity}, \code{Lanczos/Ritz.lean}:92\\
\code{eq:normalized-interface} & \code{normalized_governing_succ}, \code{Lanczos/Normalized/Recurrence.lean}:47\\
Normalized error comparison & \code{normalizedF_sub_identity}, \code{normalizedP_identity}, \code{Lanczos/Normalized/Recurrence.lean}:53,67\\
\code{eq:normalized-gram} & \code{normalized_gram_split_succ}, \code{normalizedRsucc_block}, \code{Lanczos/Greenbaum/Dilation.lean}:21,37\\
\code{eq:CKsy}, \code{eq:Kkp-power} & \code{greenbaum_CK_blocks}, \code{greenbaum_Kkp_power}, \code{Lanczos/Greenbaum/Dilation.lean}:61,118\\
\code{eq:CGC-identity} & \code{greenbaum_isometry_identity}, \code{Lanczos/Greenbaum/Dilation.lean}:109\\
\code{eq:L-commutator} & \code{greenbaumL_commutator}, \code{Lanczos/Greenbaum/Dilation.lean}:169\\
\code{eq:delta_as_projector} & \code{greenbaumDelta_projector}, \code{Lanczos/Greenbaum/Correction.lean}:54\\
\bottomrule
\end{tabularx}

\subhead{Reproduce the checks}
From \code{lean_audit/lean}, run:
\begin{verbatim}
lake build
lake env lean ../audit_report/AxiomAudit.lean
\end{verbatim}
From \code{lean_audit}, run:
\begin{verbatim}
python3 audit_report/check_zero_ritz.py
pdflatex -interaction=nonstopmode -halt-on-error \
  -output-directory=audit_report audit_report/statement_audit.tex
\end{verbatim}
The report is standalone; it does not require the paper's missing preamble,
figures, bibliography, or \code{AI} include.

\subhead{Evidence retained alongside this PDF}
\code{build.log}, \code{axiom_audit.log}, \code{AxiomAudit.lean},
\code{source_scan.txt}, \code{source_hashes.json}, \code{import_audit.json},
\code{check_zero_ritz.py}, and \code{zero_ritz_check.log}.
The axiom audit uses Lean's \code{collectAxioms} on every imported declaration
with the namespace prefix; private dependencies are covered transitively when
used by those declarations. It rejects any axiom outside the stated allowlist.

\subhead{Snapshot identification and limits}
Toolchain: \code{leanprover/lean4:v4.32.2}. Mathlib manifest revision:
\code{905b95818eb32af7874a58b427f50c1711a5e96c}.
The dependency packages are provided by the existing local symlink/cache;
this was not a clean bootstrap or an independent rebuild of Mathlib.
Paper SHA-256:
\code{18244124f5f6aa17a9324ea298a75aa1959136a390404dc169b9a845d9344e28}.

The source manifest records individual hashes for the paper and all 42 Lean
files. The prose-to-formula comparison is a human-readable semantic audit,
not a mechanized translation of LaTeX into Lean. Kernel acceptance of Lean
proofs and correspondence of their statements to the paper are separate checks;
this report records both and states the remaining gaps explicitly.

\end{document}
